% PAQUETES
%%%%%%%%%%%%%%%%%%%%%%%%%%%%%%%%%%%%%%%%%%%%%%%%%%%%%%%%%%%%%%%%%
\usepackage[utf8]{inputenc}
\usepackage[T1]{fontenc}
\usepackage[spanish]{babel}
\addto\captionsspanish{
    \renewcommand{\refname}{Bibliografía}
    \renewcommand{\tablename}{Tabla}
    \renewcommand{\listtablename}{Índice de tablas}
    \renewcommand{\listfigurename}{Índice de figuras}
    \renewcommand{\contentsname}{Índice}
}
% TikZ usa > y < en flechas (>=, ->, ++(...)); babel-español los activa como atajos.
\AtBeginDocument{\shorthandoff{>}\shorthandoff{<}}
\usepackage[margin=2cm]{geometry}
\usepackage{microtype}
% hyperref va al final del preámbulo (después de tocloft/titlesec) para no romper el .toc
\usepackage{graphicx}
% Imágenes: por defecto buscar primero en memoria/figures/, luego en memoria/
\graphicspath{{figures/}{./}}
\usepackage{booktabs}
\usepackage{enumitem}
\usepackage{array}
\usepackage{xcolor}
% parskip parchea \@starttoc; debe cargarse DESPUÉS de tocloft (ver aviso tocloft "bailing out")
\usepackage{float}
\usepackage{tabularx}
\usepackage{ltablex}
\usepackage{colortbl}
% pdflscape retirado: la versión actual de TinyTeX requiere primitivas PDF
% (\IfPDFManagementActiveF) no disponibles en este formato; no se usa
% ningún entorno landscape en el cuerpo. Si lo necesitas, usa \usepackage{lscape}.
\usepackage{needspace}
\usepackage{tikz}
\usepackage{pgf-pie}
% Evita conflictos de babel (p. ej. español: > activo) con TikZ (>=, ->, etc.)
\usetikzlibrary{babel}
\usetikzlibrary{positioning,arrows.meta,fit,backgrounds,calc,shapes.geometric,decorations.pathreplacing}
\usepackage{eurosym}
\usepackage{amsmath}
\usepackage{amssymb}
\usepackage{titlesec}
\usepackage{tocloft}
% \keepXaliases (ltablex) no está definido en todas las versiones del paquete
\makeatletter
\@ifundefined{keepXaliases}{}{\keepXaliases}%
\makeatother

\usepackage{parskip}
\setlength{\parskip}{7pt plus 2pt minus 1pt}

% Espaciado de títulos
\titlespacing*{\section}{0pt}{10pt}{4pt}
\titlespacing*{\subsection}{0pt}{8pt}{3pt}
\titlespacing*{\subsubsection}{0pt}{6pt}{2pt}
\titleformat{\section}{\normalfont\large\bfseries}{\thesection.}{0.4em}{}
\titleformat{\subsection}{\normalfont\normalsize\bfseries}{\thesubsection.}{0.4em}{}
\titleformat{\subsubsection}{\normalfont\small\bfseries}{\thesubsubsection.}{0.4em}{}

% Anchos del índice (tocloft debe cargarse antes que hyperref)
\setlength{\cftsecnumwidth}{2.8em}
\setlength{\cftsubsecnumwidth}{3.8em}
\setlength{\cftsubsubsecnumwidth}{4.8em}
\setlength{\cftfignumwidth}{2.8em}
\setlength{\cfttabnumwidth}{2.8em}
\setlength{\cftbeforesecskip}{1pt}

\setcounter{secnumdepth}{2}   % numerar solo section y subsection
\setcounter{tocdepth}{2}      % mostrar en el índice solo hasta subsection

\widowpenalty=10000
\clubpenalty=10000

% Colores reutilizados en diagramas
\definecolor{bcteacher}{RGB}{220,90,90}
\definecolor{bcstudent}{RGB}{90,140,220}
\definecolor{bcrouter}{RGB}{240,170,70}
\definecolor{bcdata}{RGB}{90,180,140}
\definecolor{bcmeasure}{RGB}{160,120,200}
\definecolor{bcneutro}{RGB}{210,210,215}

% Enlaces PDF (último paquete relevante del preámbulo)
\usepackage{hyperref}
\hypersetup{
    colorlinks=true,
    linkcolor=black,
    citecolor=black,
    urlcolor=blue
}
\usepackage{bookmark}
